\section{Implicit libraries and a sufficient repair potential}
\label{sec:types-energy}

\subsection{Exact participation types}

\Cref{thm:fpt} concerns a fixed selected named set.  To avoid conflating that
problem with a search through a large library, fix the following separate
parameterized problem.  The input is a finite target \(T\), a finite named
cell-list library \(\mathcal Q\), the explicit orientation action, an integer
\(Q\), and one of the three predicates in \cref{def:predicates}.  The question
is whether there is an area-compatible selected set \(E\subseteq\mathcal Q\)
whose canonical largest name \(L(E)\) satisfies
\[
 |T|-|L(E)|\leq Q
\]
and whose margin fiber has the requested predicate.  The classification
version returns one compressed record per behavior class; it does not list
all named selected sets.

Fix a candidate largest name \(L\), order its legal poses as
\(\mathcal A_L=(A_1,\ldots,A_s)\), and put \(q_L=|T|-|L|\).  A residual name
\(P\) is \emph{eligible for \(L\)} when \(|P|<|L|\), or when
\(|P|=|L|\) and the fixed tie order chooses \(L\) before \(P\).  For
\(A_j\in\mathcal A_L\), define
\[
 b_j=X(T)-X(A_j),\qquad
 \mathcal M_j(P)=\{B\in\PP_P:X(B)\leq b_j\text{ coordinatewise}\}.
\]
Every mask in \(\mathcal M_j(P)\) lies in the active carrier \(G_{A_j}\).
Use the common labelled carrier
\[
 S_L=\bigcup_{j=1}^{s}G_{A_j},
 \qquad |S_L|\leq h(Q+1)Q^2.
\]

\begin{definition}[Exact participation type]\label{def:participation-type}
The exact participation type of a residual name \(P\), relative to \(L\), is
\[
 \tau_L(P)=
 \bigl(|P|,\text{eligibility},
       (\mathcal M_1(P),\ldots,\mathcal M_s(P))\bigr).
\]
Each placement is stored as its actual labelled cell mask in \(S_L\), not as
an index local to one scenario.  Thus the record retains the identity of a
placement mask that occurs for more than one largest-piece pose.  Let
\(m_L(\tau)\) be the number of available residual names of type \(\tau\).
\end{definition}

\begin{lemma}[Type-isomorphism lemma]\label{lem:type-isomorphism}
Fix \(L\).  Suppose two residual named multisets
\((P_1,\ldots,P_r)\) and \((P'_1,\ldots,P'_r)\) admit a bijection \(\sigma\)
with \(\tau_L(P_i)=\tau_L(P'_{\sigma(i)})\) for every \(i\).  Replacing each
name by its matched name and retaining the same labelled placement mask gives
an isomorphism of margin fibers.  This isomorphism maps exact covers to exact
covers and is a graph isomorphism for both \(\GG^{(2)}\) and
\(\GG^{(\leq2)}\).
\end{lemma}

\begin{proof}
For every ordered pose \(A_j\) of \(L\), equality of types gives exactly the
same allowed cell masks in every matched residual coordinate.  Hence the
coordinatewise renaming is a bijection of placement tuples and preserves the
sum of margin vectors.  Because the actual masks are retained, it also
preserves the coverage function and therefore the exact-cover sublocus.
Two tuples differ in a named coordinate before the renaming exactly when
their images differ in the matched coordinate.  This remains true when the
largest pose changes: cross-scenario mask identity records which unchanged
residual coordinates are literally the same placement.  Hamming distance
one or two, and therefore both replacement graphs, are preserved.
\end{proof}

\begin{proposition}[Compressed finite-library decision]\label{prop:types}
The existential problem above and its compressed behavior classification are
FPT in \((Q,h)\).  For a candidate \(L\), a compressed record contains \(L\),
an exact participation-type count vector \(k=(k_\tau)_\tau\), the number
\[
 \prod_\tau \binom{m_L(\tau)}{k_\tau}
\]
of named residual selections represented by that vector, and the resulting
fiber predicates and counts.
\end{proposition}

\begin{proof}
For each candidate \(L\) with \(0\leq q_L\leq Q\), the ordered list of
largest poses and the carrier \(S_L\) have size bounded by a function of
\((Q,h)\).  Area, eligibility, and the exact mask family for every ordered
largest pose therefore yield only \(f(Q,h)\) possible types.  The selected
residual names are nonempty and have total area \(q_L\), so at most \(Q\) of
them are chosen.  We may truncate the available multiplicity of every type
at \(Q\), enumerate the type-count vectors satisfying
\[
 \sum_\tau k_\tau\,\operatorname{area}(\tau)=q_L
\]
and the eligibility rule, and run \cref{thm:fpt} on one representative of
each vector.  The type-isomorphism lemma proves that this representative has
exactly the same fiber, exact-cover sublocus, and move graphs as every named
selection counted by the displayed product.  Constructing the signatures
for all candidate largest names is polynomial in the explicit input size;
all superpolynomial enumeration is confined to a function of \((Q,h)\).
\end{proof}

\begin{example}[Why literal listing is output-sensitive]
Take a \(1\times2Q\) target, one named \(Q\)-bar, and \(N\) named monominoes.
There are \(\binom NQ\) distinct valid named selections.  Consequently the
literal list of all selections cannot, in general, be produced in total FPT
time independent of output size.  No Delay-FPT enumerator is claimed here;
the distinction follows the standard parameterized-enumeration vocabulary
\cite{CreignouEtAl2017}.
\end{example}
\subsection{Overlap energy}

For \(x\in\FFib_E\), define
\[
 \Phi(x)=\sum_{u\in T}\max(c_x(u)-1,0).
\]

\begin{proposition}[Energy identity and strict descent]\label{prop:energy}
For every margin node \(x\),
\[
 \Phi(x)
 =|T|-\left|\bigcup_i A_i\right|
 =\frac12\|c_x-\one_T\|_1,
 \qquad 0\leq\Phi(x)\leq q.
\]
Moreover, \(\Phi(x)=0\) exactly when \(x\in\TT_E\).  If every non-tiling node
has an exact-two neighbor of strictly smaller energy, then the fiber is
repair-exact and \(\rho(E)\leq q\).
\end{proposition}

\begin{proof}
The chosen pose areas sum to \(|T|\), so
\[
 \sum_u c_x(u)=|T|.
\]
The overlap mass \(\sum_u\max(c_x(u)-1,0)\) is therefore the total area minus
the size of the covered union.  Also
\(\sum_u(c_x(u)-1)=0\): positive deviation equals uncovered mass, proving the
\(\ell_1\) identity.  The largest pose covers \(|T|-q\) distinct cells, so at
most \(q\) cells can be missing from the union, whence \(0\leq\Phi\leq q\).
Energy zero is equivalent to unit coverage everywhere.  Finally, repeated
strict descent in the nonnegative integer \(\Phi\) reaches zero in at most
\(\Phi(x)\leq q\) steps.
\end{proof}

Strict descent is sufficient, not necessary.  Certified P21 and P26 rows
contain repairable nonzero local minima that require a nondecreasing or
uphill step.  Thus local minima do not characterize trapped components.
